\section{The signature determines the incidence coefficient}
\label{sig:section}

The incidence construction has reduced the bundle calculation to
one odd integer $s$. The signature theorem determines it in every
quaternionic rank. The small ranks follow from exact coefficients;
a single contour estimate treats all larger ranks.

\begin{theorem}\label{sig:coefficient}
In Proposition~\ref{inc:normal}, $|s|=1$. In particular, the
actual tautological four-plane bundle of every tangent great-sphere
fibration at a point of $M$ satisfies $p_1=\pm2e$.
\end{theorem}

\begin{proof}
Keep the choices $s\ge1$ and $\int_Mu^r=1$. Set
\[
 a_+a_-=1,\qquad a_+^2+a_-^2=2s,\qquad
 a_+\ge a_->0,\qquad c=2\sqrt{s},
 \quad f(y)=\frac{\sqrt y}{\tanh\sqrt y}.
\]
Hirzebruch's signature theorem
\citep[Signature Theorem~19.4]{MilnorStasheff1974}, applied to
Theorem~\ref{inc:pontryagin} gives
\begin{equation}\label{sig:formula}
 \sigma_r=[u^r]
 \frac{(f(a_+^2u)f(a_-^2u))^{r+1}}{f(c^2u)}.
\end{equation}
Here $f$ is its power series at zero; the two factors reflect
$1+2su+u^2=(1+a_+^2u)(1+a_-^2u)$.
The integral cohomology ring of $M$ gives
\[
 \sigma_r=\begin{cases}1,&r\text{ even},\\0,&r\text{ odd}.
 \end{cases}
\]
Indeed for even $r$ the middle pairing has the positive basis
$u^{r/2}$, and for odd $r$ its middle cohomology is zero.
Expanding $f(y)=1+y/3-y^2/45+2y^3/945+\cdots$ in
\eqref{sig:formula} gives
\[
 \sigma_2-1=\frac{8(s^2-1)}{15},\qquad
 \sigma_3=-\frac{32s(s^2-1)}{105}.
\]
These prove the assertion for $r=2,3$.

We give uniform estimates for all $r\ge4$. Since $s$ is odd,
the contrary case is $s\ge3$. We express the signature as a residue
and show that the two nearest nonzero poles give a main term that
dominates the remaining contour integral, forcing
$(-1)^r\sigma_r>1$. Put
\[
 \tau=\frac{a_-}{a_+}=s-\sqrt{s^2-1},\qquad
 b=\frac c{a_+}=\sqrt{2(1+\tau^2)},\qquad
 Q(z)=\tan z\tan(\tau z),\qquad \rho=\frac\pi{2b}.
\]
Then $0<\tau\le3-\sqrt8<0.172$ and
$\sqrt2<b<1.435$.
Substitute $u=-t^2$ in \eqref{sig:formula}. Since
$f(-a^2t^2)=at/\tan(at)$ and $a_+a_-=1$, taking the
coefficient by a residue and then putting $z=a_+t$ gives
\begin{equation}\label{sig:residue}
 (-1)^r\sigma_r=
 \frac\tau b\operatorname*{Res}_{z=0}
 \frac{\tan(bz)}{Q(z)^{r+1}}\,dz.
\end{equation}
The change of variables contributes the factor
$1/(ca_+)=\tau/b$; this fixes the normalization of the residue.

\paragraph{Contour representation.}
Take the rectangle with vertical sides $\Re z=\pm X$,
$X=3\pi/4$, and horizontal sides $\Im z=\pm4$.
Apart from zero, the only poles of the integrand in this rectangle
are at $z=\pm\rho$. Zeros of $\tan z$ occur only at multiples
of $\pi$, zeros of $\tan(\tau z)$ only at multiples of
$\pi/\tau$, and the next poles of $\tan(bz)$ are outside
because $3\pi/(2b)>X$. The poles of $\tan z$ at $\pm\pi/2$
give zeros of the reciprocal denominator, not poles.
Both residues at $\pm\rho$ equal
$-b^{-1}Q(\rho)^{-r-1}$, with $Q(\rho)>0$.
The residue theorem therefore expresses \eqref{sig:residue} as
\begin{equation}\label{sig:main-term}
 (-1)^r\sigma_r=L+\mathcal E,\qquad
 L=\frac{2\tau}{b^2}Q(\rho)^{-r-1}.
\end{equation}
Symmetry of the four sides and the triangle inequality bound
$|\mathcal E|/L$ by the sum of
\begin{align}
 E_v&=\frac b\pi\int_0^4
  |\tan(b(X+iy))|
  \left(\frac{Q(\rho)}{|Q(X+iy)|}\right)^{r+1}\,dy,
  \label{sig:vertical}\\
 E_h&=\frac b\pi\int_0^X
  |\tan(b(x+4i))|
  \left(\frac{Q(\rho)}{|Q(x+4i)|}\right)^{r+1}\,dx.
  \label{sig:horizontal}
\end{align}

\paragraph{Uniform bounds.}
The same contour works for every $r\ge4$ once both boundary ratios
are bounded by constants below one: their fifth powers then control
every exponent $r+1\ge5$. We retain the $\tau$ dependence in the
numerator and denominator so these bounds remain uniform as
$\tau\downarrow0$. All finite decimals below
denote the exact indicated rational numbers. Using
$3.14<\pi<3.142$ gives
\[
 \rho<1.111,\qquad \tan\rho<2.022,\qquad
 \tau\rho<0.192,\qquad Q(\rho)<2.30\tau.
\]
For the last inequality, the increase of $\tan t/t$ and the
Taylor bounds on $[0,0.192]$ give $\tan t/t<1.013$,
and $2.022\cdot1.111\cdot1.013<2.30$.
The same monotonicity gives
\begin{equation}\label{sig:ratio}
 \frac{Q(\rho)}{\tan(\tau X)}
 \le\frac{\rho\tan\rho}{X}<0.955.
\end{equation}
On the vertical sides, $|\tan(X+iy)|=1$.
Also $\cos(2bX)\ge0$, since
$2bX=3b\pi/2\in(2\pi,5\pi/2)$, and hence
$|\tan(b(X+iy))|\le1$.
The elementary complex tangent identity gives
\begin{align*}
 \frac{|\tan(\tau(X+iy))|^2}{\tan^2(\tau X)}
 &=\frac{1+\sinh^2(\tau y)/\sin^2(\tau X)}
         {1+\sinh^2(\tau y)/\cos^2(\tau X)}\\
 &\ge 1+0.08y^2\qquad(0\le y\le4).
\end{align*}
For this last bound, the numerator is at least $1+y^2/X^2$,
whereas $\cos(\tau X)>0.917$ and
$\sinh(\tau y)\le1.082\tau y$ bound the denominator by
$1+Ay^2$ with $A=(1.082\cdot0.172/0.917)^2$.
The inequality follows from
$X^{-2}\ge A+0.08+16(0.08)A$.
Since $r+1\ge5$, \eqref{sig:ratio} now implies
\begin{align*}
 E_v
 &<\frac{1.435}{3.14}(0.955)^5
      \int_0^4(1+0.08y^2)^{-5/2}\,dy\\
 &<0.783.
\end{align*}
Here the integral equals
\[
 \frac{4(1+2\cdot0.08\cdot16/3)}{(1+0.08\cdot16)^{3/2}}
 <2.154.
\]

On a horizontal side the identity for the modulus of tangent
gives
\[
 |Q(x+4i)|\ge\tanh4\,\tanh(4\tau),\qquad
 |\tan(b(x+4i))|\le\coth(4b)<1.001.
\]
Since $\tau/\tanh(4\tau)$ is increasing,
\[
 \frac{Q(\rho)}{\tanh4\,\tanh(4\tau)}
 <\frac{2.30\cdot0.172}{0.999\cdot0.595}.
\]
Substitution in \eqref{sig:horizontal}, with $X/\pi=3/4$,
therefore gives
\[
 E_h<\frac{3\cdot1.435}{4}\,1.001
       \left(\frac{2.30\cdot0.172}{0.999\cdot0.595}\right)^5
 <0.142.
\]
The numerical trigonometric inequalities used here can be checked
without numerical approximation: use
\[
 \sin t\le t-t^3/6+t^5/120,\qquad
 \cos t\ge1-t^2/2+t^4/24-t^6/720
 \quad(0\le t\le1.111)
\]
for $\tan\rho$ and the smaller real arguments. For the
hyperbolic inequalities use
\[
 \frac{\sinh t}{t}\le1+\frac{t^2}{6}
       +\frac{t^4/120}{1-t^2/42}\quad(0\le t\le0.688).
\]
The first seven terms of the exponential series at $1.376$
give $\tanh(0.688)>0.595$, while $e^8>2000$ gives
$\tanh4>0.999$. The same bound and $b>\sqrt2$ imply
$\coth(4b)<1.001$.

Finally $Q(\rho)<2.30\tau<1$, so \eqref{sig:main-term}
and $r+1\ge5$ imply
\[
 L>\frac{\tau}{1+\tau^2}(2.30\tau)^{-5}
 \ge\frac{1}{(1+0.172^2)\,2.30^5\,0.172^4}>17.
\]
Combining these estimates yields
\[
 (-1)^r\sigma_r>(1-0.783-0.142)17=1.275>1,
\]
contradicting either possible cohomological signature. Thus $s=1$.
Restricting \eqref{inc:normal-classes} to a fixed-$p$ fiber
gives $e(\mathcal V)=-u_2$ and $p_1(\mathcal V)=2u_2$ there.
Proposition~\ref{inc:normal} identifies this restriction with
the actual tautological four-plane bundle, which proves the
last assertion as well.
\end{proof}
